% Chapter 2: Background. Draft for the lead, 23 Nov 2026. Every number is traced in ch02_background_NOTES.md.
% Citation keys: refs_extra.bib, ch02_background.bib, ch03_related.bib. Placeholders: \CITE, \FIG (defined in main.tex).
\chapter{Background: EL imaging, defect classes and crack label physics}\label{ch:background}

This chapter explains what an electroluminescence (EL) image shows (Section~\ref{sec:bg_el}), describes the BenchmarkEL
images, classes and sources (Section~\ref{sec:bg_classes}), and shows why crack and inactive labels behave so differently
(Sections~\ref{sec:bg_crack} and~\ref{sec:bg_inactive}). It then defines the metrics (Section~\ref{sec:bg_metrics}) and
links defect maps to module power (Section~\ref{sec:bg_severity}). Each measurement of our own names its label set and
cells; Chapter~\ref{ch:data} gives the full audit.

\section{Electroluminescence imaging of silicon solar cells}\label{sec:bg_el}

EL images show defects that are invisible to the eye; factories use them before a module ships, and labs and plants over
its life \citep{pratt2023benchmark}. A crystalline silicon module is a laminate of cells wired in series. Metal fingers on each cell collect the
current, busbars on the cell gather it, and soldered interconnect ribbons carry it to the next cell
\citep[Sec.~4.6]{kontges2014iea}. Cells are either mono-crystalline, with
rounded corners and no visible grain boundaries, or multi-crystalline, with square corners and grain boundaries that appear
as small dark regions \citep{pratt2023benchmark}.

In an EL measurement the module is connected to a power supply and driven in forward bias. Electrons and holes that
recombine radiatively emit photons near \SI{1100}{\nano\metre}, outside the visible range, and a camera optimised for
that band records them \citep{pratt2023benchmark}. The local emission intensity follows the minority
carrier diffusion length of the silicon \citep{fuyuki2005el}, and the highest EL signal of each cell in a module is
proportional to its operating voltage \citep{potthoff2010voltage}. Bright regions are therefore connected silicon of good
quality. Dark regions are
of three kinds: regions that the current does not reach \citep{pratt2025ictas}, regions where carriers recombine without
emitting light \citep{fuyuki2005el}, and regions covered by metal, such as the ribbons, which appear as dark lines in our
images.

Defects follow from this reading. A crack appears as a dark line. When a crack separates a part of the cell that has no
remaining connection to a ribbon, that part no longer contributes to the module power \citep{haase2018crackwidth}; it
receives no current and appears as a dark region, the inactive area. Interrupted fingers appear as short dark
streaks; BenchmarkEL labels them as gridline defects \citep{pratt2023benchmark}.

The physical crack is far narrower than anything the image resolves. Measured by digital image correlation in a standard
module under loads comparable to the \SI{5400}{\pascal} test of IEC~61215, a single crack was \SI{3.4}{\micro\metre} wide
\citep{haase2018crackwidth}, while a whole cell spans about 250~px in our images (Chapter~\ref{ch:data}). The camera records
a dark line a few pixels wide around the crack (Section~\ref{sec:bg_crack}), and the width of a crack label is a drawing
convention, not a measurement.

\section{The BenchmarkEL images, classes and sources}\label{sec:bg_classes}

BenchmarkEL \citep{pratt2023benchmark} is the main labelled set of this work. Its release \texttt{dataset\_20211104} holds
582 cell images, split by its authors into 478 train, 54 validation and 50 test cells, with a pixel label for each of
24 classes. A PV expert chose them from more than 80,000 images of five sources as a representative sample of mono- and
multi-crystalline cells with common features and defects \citep{pratt2023benchmark}. Later releases and our corrected labels
are described in Chapters~\ref{ch:related} and~\ref{ch:data}.

Each image is 512 $\times$ 512~px and shows a full cell in the centre, surrounded by parts of the adjacent cells, the
module edge or padding \citep{pratt2023benchmark}. For the three sources that supplied module images, the centre cell covers
25 to 30\% of the non-padding pixels and the labels cover the whole window; the two public sources supplied single cells in
about 75\% padding (Chapter~\ref{ch:data}). An image is thus a window of about 3 $\times$ 3 cells, and, following the
protocol (Chapter~\ref{ch:protocol}), a ``cell'' in this thesis means such a window.

\begin{table}[htbp]
\centering
\small
\caption{The 24 BenchmarkEL classes (id and name as in the official class list of release \texttt{dataset\_20211104}).
Features are components of a module; defects are extrinsic faults \citep{pratt2023benchmark}. The two focus classes of
this work are in bold.}
\label{tab:bg_classes}
\begin{tabular}{>{\raggedright}p{2.0cm} >{\raggedright\arraybackslash}p{11.5cm}}
\toprule
Kind & Classes (id name) \\
\midrule
Features (12) & 0 background, 1 sp multi, 2 sp mono, 3 sp dogbone, 4 ribbons, 5 border, 6 text, 7 padding, 8 clamp,
9 busbars, 21 frame edge, 22 jbox \\
Defects (12) & 10 crack rbn edge, \textbf{11 inactive}, 12 rings, 13 material, \textbf{14 crack}, 15 gridline, 16 splice,
17 dead cell, 18 corrosion, 19 belt mark, 20 edge dark, 23 meas artifact \\
\bottomrule
\end{tabular}
\end{table}

Table~\ref{tab:bg_classes} lists the classes; a mask stores the class id of every pixel. On the release masks, background and padding together take 86.5\% of the train pixels, four classes
(clamp, crack rbn edge, rings and meas artifact) have under 1,300 pixels in all 478 train images, and 7 classes are absent
from the validation cells and 8 from the test cells. On the release masks crack takes 0.19\% of the train pixels; on the revised labels (label set v1,
Chapter~\ref{ch:data}), crack occurs in 209 of the 478 train cells; inactive occurs in 42 of them. An all-class mean
IoU is thus a mean over a different set of classes on each cell set, which is why Section~\ref{sec:bg_metrics} fixes the
class set.

The images come from three private sources, the CSIR Energy Centre (a South African module test laboratory), CFV Labs (a
US test laboratory) and ARTsolar (a South African manufacturer), and two public ones, the CWRU SDLE collection and ZAE
Bayern \citep{pratt2023benchmark}; the release holds ARTS 195, CFVS 187, CSIR 122, BMRK 60 and SDLE 18 images (Figure~\ref{fig:F1} shows one pool cell of each). The BMRK images are cells of the
public ELPV dataset \citep{deitsch2019elpv}: each is the ELPV cell with the same number, rotated by 90 degrees and downscaled
from 300 to 256~px inside the padded canvas (four of four cells checked, normalised cross-correlation 1.000).

The sources differ in more than origin. CSIR images have the lowest contrast (median intensity standard deviation 24,
against 34 to 48 for the other sources), crack labels have source-specific widths (Section~\ref{sec:bg_crack}), and on
the training pool 10 of the 24 classes occur in one source only. A classifier on global image statistics names the source
with balanced accuracy 0.94 against a chance level of 0.20 (Table~\ref{tab:T19a}), so a model can recognise the source,
and with it the labelling convention, without looking at the defect.

\FIG{F1}{One pool cell per source (ARTS, BMRK, CFVS, CSIR, SDLE), chosen by a fixed rule (median crack size) and not by eye, with the
v1.1 crack (vermillion, drawn 1~px dilated) and inactive (blue) labels. ARTS, CFVS and CSIR images are windows of about 3 $\times$ 3
cells; BMRK and SDLE are single cells on white padding. The release's one double-annotated image is excluded.}

\section{Crack labels: physics, convention and geometry}\label{sec:bg_crack}

Three widths meet in a crack label: the physical crack, micrometres wide (Section~\ref{sec:bg_el}); the dark line in the
image; and the line the annotator drew. On the release masks, measured at 9,302 points of the label skeletons, the visible
dark line has a median full width at half maximum of 4.0~px (interquartile range 3.0 to 4.75), while the label has a median
width of 1.75~px (interquartile range 1.25 to 2.5). At 77\% of the points the visible line is at least 1~px wider than its
label. The annotators traced a line of 1 to 2~px somewhere on a band about 4~px wide.

Where on the band the line was drawn varies. At 61,176 label skeleton points, the darkest pixel along the normal lies within
$\pm$1~px of the label at 48\% of the points (ARTS 75\%, BMRK 38\%) and within $\pm$2~px at 64\%, against 11\% and 21\%
for random background points: the labels are on the cracks, but half of the label points are more than 1~px off the
darkest ridge.

The label width is a convention of each source: the median label width on the release masks is 1.5~px for ARTS, 1.75~px
for CFVS and CSIR, 3.0~px for BMRK and 3.25~px for SDLE. On the training pool under label set v1.1, the audit finds the
same order with its own width measure (pooled, ARTS 1.59 to BMRK 3.09~px per centreline pixel), and the source explains 29.6\% of
the variance of the per-cell width (Table~\ref{tab:T19a}).

Thin labels make strict pixel IoU fragile. Take a perfect crack mask, move it, and score it against the original. On the
legacy test masks (release labels), a 1~px shift gives crack IoU 0.559 (mean of four directions), a diagonal 1~px shift
0.434 and a 2~px shift 0.336. On the training pool under v1.1, pooled
over its 229 crack cells, the 1~px shift gives 0.483 and the 2~px shift 0.249; the same label reduced to its 1~px
centreline scores 0.465, and redrawn 3~px wide 0.605. The tolerance metric relF1@2 of Section~\ref{sec:bg_metrics} is 1.0
for every one of these variants (Figure~\ref{fig:F2}). A label that is correct to within one pixel thus loses about half of its
strict IoU.

Annotators also disagree with each other. The two image pairs of the release whose copies were annotated independently
reach crack IoU 0.257 and 0.210 but tolerance F1 at 2~px of 0.82 and 0.99; with two pairs this is an illustration, not an
estimate. On the training pool under v1.1, the one double annotation agrees with itself at strict
crack IoU 0.259 and relF1@2 0.829 (Table~\ref{tab:T19a}).

Most cracks are small and some are faint. On the release masks, 76\% of the crack components in the train cells, 59\% in the
validation cells and 88\% in the test cells are under 20~px, and 7.2\% of all crack components (16\% for CSIR) are less
than 2\% darker than their surroundings. Some continuous cracks are labelled as dashed traces, so a model that predicts
the whole crack is scored wrong on every gap, and some visible cracks are not labelled at all. Chapter~\ref{ch:data}
counts such label-error candidates on the pool.

These properties shape what a crack score measures. For the best model of the pre-protocol work, on the 49 clean legacy test
cells (in-distribution, known leakage, context only), 70\% of the strict crack IoU loss, 0.358 of 0.508, came from pixels
within 2~px of the ground truth; missed cracks made up 18\% of the loss and spurious predictions 11\%. PF's out-of-fold pool
predictions (the baseline of Chapter~\ref{ch:methods}, not the model above) score strict crack IoU 0.233, 0.365 and 0.414 when the labels are redrawn as a 1~px centreline, kept as v1.1, or
drawn 3~px wide, while relF1@2 moves only from 0.824 to 0.849 (Table~\ref{tab:T19a}). We read this as follows: strict crack
IoU mixes detection with the line-width convention of the annotators, and a distance-tolerant metric separates the two
better. Chapter~\ref{ch:protocol} turns this into the rule that relF1@2 decides every crack contrast and that strict IoU is
printed with its convention band.

\FIG{F2}{Crack label geometry on the training pool (v1.1, 229 crack cells): strict crack IoU of the label against itself shifted by
1~px and by 2~px (mean of 8 directions), reduced to its 1~px centreline and redrawn 3~px wide. Bars are pooled values, markers per
source (none for the 2~px shift). relF1@2 is 1.0 for every variant.}

\section{Inactive area}\label{sec:bg_inactive}

An inactive area is a part of a cell that a crack has cut off electrically. It produces no power
\citep{haase2018crackwidth}, and regions darkened by partly or fully disconnecting cracks are the ones that go with module
power loss \citep{pratt2025ictas}. Of all classes, it is the most directly tied to the electrical consequence of a defect.

Its labels behave the opposite way to crack labels. On the release masks, inactive never occurs in an image without a crack
label, in any split, and 99\% of the inactive components touch a crack label within 3~px (440 of 445 in the train cells,
18 of 18 in the validation cells, 51 of 52 in the test cells): inactive regions are bounded by cracks. 20\% of the crack
pixels of the train cells lie in or at inactive regions, so a single-label mask must give such pixels one of the two classes.
The inactive labels are consistent where crack labels are not: on the legacy test masks, a perfect inactive mask shifted by
1~px keeps IoU 0.943, and the duplicate pairs agree at inactive IoU 0.64 with tolerance F1 0.92 at 2~px.

The difficulty is scarcity. Inactive occurs in 42 of the 478 train cells, 6 of the 54 validation cells and 16 of the 50 test
cells (Table~\ref{tab:T3b}), and on the test cells three cells hold 59\% of the inactive pixels. Any inactive score on a
legacy set is therefore decided by a handful of cells. In our reading, crack is limited by label noise and inactive by the
number of examples; Chapters~\ref{ch:protocol} and~\ref{ch:methods} treat the two classes with different rules for this
reason.

\section{Metrics}\label{sec:bg_metrics}

The metrics of the protocol are computed at a fixed operating point (argmax unless stated) from counts pooled over a cell
set. Pixels marked ignore in the label set are not scored.

\paragraph{IoU and mean IoU.} For a class $c$, one confusion matrix over all scored pixels of the cell set gives the true
positives $\mathit{TP}_c$, false positives $\mathit{FP}_c$ and false negatives $\mathit{FN}_c$, and
\begin{equation}
  \mathrm{IoU}_c = \frac{\mathit{TP}_c}{\mathit{TP}_c + \mathit{FP}_c + \mathit{FN}_c}.
\end{equation}
We call crack $\mathrm{IoU}_c$ \emph{strict} crack IoU, to set it apart from the tolerance metrics below. A mean IoU is only
defined together with its class set. A class without ground truth in the scored set is never scored as 0 or 1; an all-class
mean is reported only as the mean over the $k$ classes with ground truth, with $k$. The balance metric of this thesis,
mIoU-R13, is the mean IoU over 13 fixed classes: background, ribbons, padding, border, gridline, sp mono, crack, sp multi,
material, frame edge, splice, inactive and edge dark. These are the classes that have ground truth in at least three of the
legacy test cells.

\paragraph{Tolerance metrics for crack.} Let $G$ be the set of ground-truth crack pixels and $P$ the set of predicted crack
pixels, and let $d(p, S)$ be the Euclidean distance from pixel $p$ to the nearest pixel of $S$ in the same image (computed by
a distance transform per image). With tolerance $r$ in pixels,
\begin{equation}
\begin{gathered}
  \mathrm{relP@}r = \frac{\lvert\{p \in P : d(p, G) \le r\}\rvert}{\lvert P \rvert}, \qquad
  \mathrm{relR@}r = \frac{\lvert\{g \in G : d(g, P) \le r\}\rvert}{\lvert G \rvert}, \\
  \mathrm{relF1@}r = \frac{2\,\mathrm{relP@}r \cdot \mathrm{relR@}r}{\mathrm{relP@}r + \mathrm{relR@}r},
\end{gathered}
\end{equation}
with the counts in each numerator and denominator summed over all cells of the set. Recall is never reported without
precision. The tolerance of 2~px matches the 95th percentile of the distance
between two drawings of one crack on the few such pairs we have (Chapter~\ref{ch:protocol}). A prediction drawn somewhat
wider than the label stays within the tolerance, so these metrics do not see line width; the protocol prints the ratio of
predicted to labelled crack pixels beside them as a width guard.

Published EL papers aggregate differently: per image with the median over images, pooled with or without background, and
with different rules for absent classes. Chapter~\ref{ch:related} shows how far this moves the numbers.

\section{From defect maps to severity}\label{sec:bg_severity}

A crack matters for module power mainly through the area it isolates. Köntges et al.\ grouped cell cracks by whether they
separate parts of a cell: mode~A cracks isolate nothing, mode~B cracks partly disconnect a region and mode~C cracks
disconnect it fully; in laboratory stress tests, power loss grew as mode~A cracks turned into modes B and C
\citep{kontges2011risk}, as summarised by \citet{pratt2025ictas}. Greater power loss goes with more dark area in the EL image,
including regions darkened by mode~B and mode~C cracks \citep{pratt2025ictas}. Learned models predict module power from EL
images with inactive regions and cracks as the main defects \citep{hoffmann2021power}, and a crack segmentation can be
turned into a worst-case inactive area per cell \citep{chen2022pvsc}. We read the cited work as saying that, for severity,
the inactive map and whether a crack closes off an area matter more than the exact pixels of a crack line.

Thermal imaging sees the electrical consequence rather than the crack line. In tests on 12 cells, Dhimish and Lazaridis
found no hotspot on crack-free or micro-cracked cells, while cells with shaded or broken-down areas were likely to develop
hotspots, with cell temperatures rising from 25\,$^{\circ}$C to 100\,$^{\circ}$C \citep{dhimish2021hotspots}. This matters
for the industrial partner of this work, which offers drone thermal inspections and, as of November 2026, has no EL rig
(Chapter~\ref{ch:intro}).

No labelled cell in this work has an electrical measurement, so the thesis gives no watt figure or power band per cell.
Chapter~\ref{ch:deployment} reports what the segmentation output can and cannot say about severity (Table~\ref{tab:T18}).
